\documentclass[conference]{IEEEtran}
\usepackage{times}

\usepackage[numbers]{natbib}
\usepackage{multicol}
\usepackage{graphicx}
\usepackage{booktabs}
\usepackage{multirow}
\usepackage{xcolor}
\let\labelindent\relax
\usepackage{enumitem}

\usepackage[bookmarks=true]{hyperref}

\pdfinfo{
   /Author (Anonymous)
   /Title  (AutoDex Rebuttal)
   /CreationDate (D:20260327120000)
   /Subject (Robotics)
   /Keywords (Dexterous;Grasping)
}

\definecolor{R1}{RGB}{0,128,0}
\definecolor{R2}{RGB}{0,0,200}
\definecolor{R3}{RGB}{200,0,0}
\definecolor{AC}{RGB}{128,0,128}

\newcommand{\rone}{{\color{R1}R1}}
\newcommand{\rtwo}{{\color{R2}R2}}
\newcommand{\rthree}{{\color{R3}R3}}
\newcommand{\acmark}{{\color{AC}AC}}

\renewcommand{\thetable}{R-\arabic{table}}
\renewcommand{\thefigure}{R-\arabic{figure}}

\newcommand{\rssparagraph}[1]{{\bf #1.}}

% Allow up to 4 floats per column top, using up to 95% of column
\setcounter{topnumber}{4}
\renewcommand{\topfraction}{0.95}
\setcounter{totalnumber}{4}
\setlength{\belowcaptionskip}{-10pt}
\setlength{\textfloatsep}{5pt}

\begin{document}

% \title{AutoDex [Paper-ID: 99]\vspace{-1cm}}
% \maketitle
% \thispagestyle{empty}

%% ============================================================
%% TEXT BODY
%% ============================================================
% AutoDex [Paper-ID: 99]
\noindent [AutoDex, Paper-ID: 99]
We thank the reviewers for recognizing the importance of bridging the sim-to-real gap ({\color{R1}R1=zCHJ}, {\color{R3}R3=CUok}), the ``clear quantitative argument for real-world filtering'' and ``coverage-preserving selection design'' ({\color{R2}R2=5XQi}), and the inclusion of both successful and failed grasps ({\color{R3}R3}). We address all concerns below.


%% --- SECTION 2: Necessity of real-world validation ---

\noindent [\textcolor{cyan}{All}] \rssparagraph{Necessity of real-world validation}
As shown in Sec. V.B and Table II of our manuscript, incorporating physically validated data through our AutoDex pipeline substantially improves performance compared to simulation-only training (average success rate: 71\% vs. 33\%). This gap becomes even more pronounced in complex scenes.
%
To further verify the effectiveness of real-world validation, we conduct additional experiments with a larger set of objects (20 in total), as illustrated in Fig. R-1 (b) and (c). Using the same AutoDex framework, we collect an average of 50 physically validated samples per object. We analyze success rates in 515 test trials across scene constraint types as shown in Fig. R-1 (c), where the plot the performance per materials, scenes, weight variations. The results consistently demonstrate strong performance gains with AutoDex (76\% vs. 34\%), showing similar performace gap in Table II of our original manuscript.

To more clearly highlight the limitations of simulation-only approaches, we further analyze precision--recall characteristics by comparing simulation predictions with real-world success and failure outcomes. Specifically, we plot PR curves by varying simulation parameters and validity thresholds (e.g., object weight, friction, and number of trials), as shown in Fig. R-1 (b). The results reveal consistent limitations across all thresholds. For example, even when sacrificing recall, precision remains below 70\%. This indicates that domain randomization alone is insufficient to ensure real-world physical validity, leading to frequent false positives and false negatives. We will include these additional experiments and discussions in our revision.


%% --- SECTION 1: Geometric Transfer ---
\noindent [\acmark, \rtwo, \rthree] \rssparagraph{Clarification on geometric transfer}
We define geometric feasibility as satisfying inverse kinematics (IK) solvability, collision avoidance, and kinematic reachability. These properties can be reliably verified in simulation and generally transfer well to the real world, effectively filtering out clearly infeasible cases. However, passing geometric feasibility does not guarantee successful grasping in practice. This limitation arises from the imperfect modeling of contact dynamics in simulation. We will further clarify these definitions and distinctions in the revised manuscript.

\noindent [\acmark, \rthree] \rssparagraph{Why still rely on simulation}
Our simulation filter operates under deliberately optimistic settings so that it removes only obvious failures rather than acting as a precise predictor. As shown in \emph{Fig.~R-1,(b)}, simulation produces both false negatives and false positives and cannot be made fully reliable at any threshold. Our strategy is therefore two-staged: simulation narrows the candidate set by discarding clear failures, and subsequent real-world execution filters out the remaining false positives.


\begin{figure}[t]
\includegraphics[width=\columnwidth]{figures/merged.png}%
\vspace{-24pt}
\caption{(a): (L) Scenario solvability: Ours vs Baseline by grasp budget $N$, with real-world validation at $N$=10,30. (R) Constrained scene execution examples.
(b): (L) Precision--recall of sim DR; no threshold achieves both high precision and recall. (R) Automated grasp execution sequence.
(c): (L) Expanded object set (20 objects). (R) Physical validation: success rate by material, scene type, and weight ($\pm$95\% CI).
(d): (L) Collection efficiency: time \& cost vs samples (AutoDex vs Teleop). (R) Perception quality vs camera count.}
\label{fig:merged}
\end{figure}

\noindent [\acmark, \rtwo] \rssparagraph{Real-world Scene Coverage Test}
Following \rtwo's feedback, we additionally conduct
real-world experiments to validate that AutoDex improves scenario solvability, corresponding to the real-world counterpart of Tab.~I in our paper. Specifically, we perform grasping experiments using 3 objects across 3 constraint types, with 10 configurations each (90 trials in total).
We measure \emph{expected scene coverage}, defined as the probability that at least one feasible grasp, verified as successful in the real world, exists within a given budget $N$. As shown in \emph{Fig.~R-1 (a)}, we plot this metric while varying $N$ up to 100 (x-axis). Our AutoDex dataset consistently outperforms the baseline across all budgets.
We will include these additional experiments and discussions in the revised manuscript.


\noindent [\acmark, \rtwo] \rssparagraph{Efficiency compared to teleoperation}
Two methods exist for real-world grasp validation: teleoperation and our automated pipeline. AutoDex achieves [X] grasps/hr vs [X] for teleop, requiring only [X] min of human time per grasp (object placement). Teleop requires continuous operator presence (body suit, 22-DoF control), with sessions limited to [X] hours before fatigue and Xsens charging. AutoDex runs unattended overnight and scales linearly with robots without additional operators.

Beyond efficiency, AutoDex provides \emph{diversity}: teleop datasets are biased toward operator-preferred approach paths, and constructing constrained scenes (wall, shelf) during teleop is impractical. AutoDex evaluates grasps under procedural constraints via geometric transfer, and is generation-method agnostic --- any generator (BODex, RL, human priors) can be plugged in. We focus on stable grasps in this work --- an intentional scope decision, as no physically-validated dexterous grasp dataset existed even for this foundational task. The pipeline extends to functional grasps by swapping the generation method; validation and selection remain unchanged.


%% --- SECTION 5: Camera scalability ---

\noindent [\acmark, \rthree] \rssparagraph{Scalability of the multi-camera setup}
The 24-camera setup is a research choice, not a pipeline requirement. Our system is agnostic to camera count and type. Calibration: $\sim$40 board images ($\sim$10 min) + automated hand-eye ($\sim$10 min), regardless of camera count. During grasping, the hand occludes [X]\% of the object on average. More cameras mitigate this, but the trade-off is per-lab (Table~III: 4-cam ADD-S 1.17mm / 51\% success vs 24-cam 0.31mm / 62\%).

\noindent [\acmark, \rtwo, \rthree] \rssparagraph{Performance gains through physical filtering}
We expanded evaluation with more objects, more trials, and confidence intervals (Fig.~R-1 (c)). Real-filtered grasps achieve [X]\% $\pm$ [X] vs [X]\% $\pm$ [X] for sim-validated (Fisher's exact $p < $ [X]), with consistent gains across all material and scene categories. This confirms that physical filtering captures dynamics that simulation cannot predict.


\end{document}
